\subsection{不可公度尺度、临界线局域化与记录轴最小闭包的非周期刚性}\label{subsec:conclusion-aperiodic-critical-rigidity-multiscale-ledger}
本小节把双底去周期化、尺度 Mellin 叠加完成化、对数 Sobolev--Gross 超收缩证书、连续横截方向的唯一性以及扩展保存记录轴的最小乘积闭包压缩为同一条终端链。
其共同结论是：可审计谱几何既不退化为单底算术梳齿，也不允许在离临界线方向任意漂移；一旦自对偶性、混合强度与记录轴纪律同时进入闭包，稳定可见结构便呈现出一种“临界线支持而竖向非晶格”的刚性形态。

\begin{theorem}[两不可公度底摧毁竖直晶格对称]\label{thm:conclusion-two-incommensurable-bases-destroy-vertical-lattice}
取 \(L_1,L_2>1\) 满足
\[
\frac{\log L_1}{\log L_2}\notin\QQ,
\]
并设 \(g,h\) 不具有任何非平凡单位根乘法不变性。
定义
\[
f(s):=g(L_1^{-s})+h(L_2^{-s}).
\]
则 \(f\) 不存在任何非零纯虚周期；亦即，不存在 \(T\in i\RR\setminus\{0\}\) 使
\[
f(s+T)\equiv f(s).
\]
\end{theorem}

\begin{proof}
这正是定理 \ref{thm:xi-two-base-deperiodization-no-imag-period} 的直接重述。
该定理已经表明：仅凭两条对数不可公度的尺度，便足以摧毁单底制度所强制的严格竖向晶格对称。
\end{proof}
\leanverified{paper\_conclusion\_two\_incommensurable\_bases\_destroy\_vertical\_lattice}

\begin{theorem}[函数方程与临界线实值性并不强制竖向周期]\label{thm:conclusion-selfdual-scale-mellin-does-not-force-vertical-periodicity}
设 \(\mathcal X_{\Delta,w}\) 为定义 \ref{def:xi-scale-mellin-superposition} 的尺度 Mellin 叠加完成化，并满足命题 \ref{prop:xi-scale-mellin-no-nonzero-imag-period} 的非周期性假设。
则在自对偶条件下，\(\mathcal X_{\Delta,w}\) 同时满足严格函数方程与临界线实值性：
\[
\mathcal X_{\Delta,w}(s)=\mathcal X_{\Delta,w}(1-s),
\qquad
\mathcal X_{\Delta,w}\!\left(\frac12+it\right)\in\RR
\quad (t\in\RR),
\]
但仍不存在任何非零纯虚周期。
因此，“函数方程 \(+\) 临界线实值性”与“竖向算术周期性”在逻辑上彼此独立。
\end{theorem}

\begin{proof}
函数方程与临界线实值性由命题 \ref{prop:xi-scale-mellin-fe-real} 直接给出；
非零纯虚周期的排除则由命题 \ref{prop:xi-scale-mellin-no-nonzero-imag-period} 给出。
两者合并即得结论。
\end{proof}
\leanverified{paper\_conclusion\_selfdual\_scale\_mellin\_does\_not\_force\_vertical\_periodicity}

\begin{theorem}[自对偶尺度混合中的有限余维临界插值自由度]\label{thm:conclusion-selfdual-scale-mellin-finite-codim-interpolation}
\leanverified{paper\_conclusion\_selfdual\_scale\_mellin\_finite\_codim\_interpolation}
设 \(\Delta\) 为自对偶核，并对每个 \(t\in\RR\) 定义实线性泛函
\[
L_t:L^1_c((0,\infty),\RR)\to\RR,
\qquad
L_t(w):=
\mathcal X_{\Delta,w}\!\left(\frac12+it\right)
=
\int_0^\infty w(x)\,\Xi_{\Delta,e^x}\!\left(\frac12+it\right)\,\dd x.
\]
取互异实数
\[
t_1,\dots,t_m,t_\ast\in\RR,
\]
并假设
\[
L_{t_1},\dots,L_{t_m},L_{t_\ast}
\]
在 \(L^1_c((0,\infty),\RR)\) 上线性无关。则存在一个实值、紧支撑、分段常数且非零的权函数 \(w\)，使
\[
\mathcal X_{\Delta,w}\!\left(\frac12+it_j\right)=0
\qquad (1\le j\le m),
\]
并且
\[
\mathcal X_{\Delta,w}\!\left(\frac12+it_\ast\right)=1.
\]
与此同时，\(\mathcal X_{\Delta,w}\) 仍保持
\[
\mathcal X_{\Delta,w}(s)=\mathcal X_{\Delta,w}(1-s),
\qquad
\mathcal X_{\Delta,w}\!\left(\frac12+it\right)\in\RR
\quad (t\in\RR).
\]
因此，函数方程与临界线实值性之内仍保留严格的有限探针不可定性：有限多个临界线采样值并不足以锁死自对偶尺度混合族。
\end{theorem}
\begin{proof}
由线性无关性，映射
\[
\Lambda:L^1_c((0,\infty),\RR)\to \RR^{m+1},
\qquad
\Lambda(w):=
\bigl(L_{t_1}(w),\dots,L_{t_m}(w),L_{t_\ast}(w)\bigr)
\]
具有像空间维数 \(m+1\)。
因此可取 \(w_0,\dots,w_m\in L^1_c((0,\infty),\RR)\)，使得矩阵
\[
\mathcal M:=
\bigl(L_{t_i}(w_j)\bigr)_{i\in\{1,\dots,m,\ast\},\,j=0,\dots,m}
\]
可逆。
令 \(c=(c_0,\dots,c_m)^\top\) 为线性方程组
\[
\mathcal M c=(0,\dots,0,1)^\top
\]
的唯一解，并置
\[
w:=\sum_{j=0}^{m}c_jw_j.
\]
则由构造立得
\[
L_{t_j}(w)=0\qquad(1\le j\le m),
\qquad
L_{t_\ast}(w)=1,
\]
故 \(w\not\equiv 0\)。

再利用紧支撑分段常数函数在 \(L^1_c((0,\infty),\RR)\) 中的稠密性，可把每个 \(w_j\) 同时替换为足够接近的实值紧支撑分段常数函数而不破坏 \(\det(\mathcal M)\neq 0\) 这一开条件。
因此可直接取 \(w\) 也为实值、紧支撑、分段常数。

最后，自对偶完成化的函数方程与临界线实值性由命题 \ref{prop:xi-scale-mellin-fe-real} 对任意实权 \(w\) 都成立，故上述插值自由并不破坏
\[
\mathcal X_{\Delta,w}(s)=\mathcal X_{\Delta,w}(1-s),
\qquad
\mathcal X_{\Delta,w}\!\left(\frac12+it\right)\in\RR.
\]
于是有限个临界线取值条件只能切出一个有限余维仿射截面，而不能消灭整个自对偶尺度混合自由度。证毕。
\end{proof}

\leanverified{paper\_conclusion\_logsobolev\_mixing\_forces\_critical\_localization}
\begin{theorem}[混合强度的非退化强制临界线局域化]\label{thm:conclusion-logsobolev-mixing-forces-critical-localization}
在“视界零知识---幺正切片”算子模型中，设不可见部分上存在完全正、迹保持半群 \((\mathcal T_u)_{u\ge 0}\)，并满足维数无关的对数 Sobolev 不等式，其常数记为 \(\alpha>0\)。
若视界算子写为
\[
K_s\big|_{\mathfrak H_{\mathrm{inv}}}
=
\int_0^\infty e^{-(s-\frac12)u}\,\mathcal T_u(U)\,\dd u,
\]
则存在仅依赖正规化的显式常数 \(c>0\)，使得带外区域
\[
\left|\Re(s)-\frac12\right|>\frac{c}{\alpha}
\]
上无零点事件。
特别地，若沿某一尺度族保持 \(\inf \alpha>0\)，则全部稳定零点事件都被压缩进一条以 \(\Re(s)=\frac12\) 为中心的有界束缚带。
\[
\text{更精确地，若某稳定零点事件与临界线保持距离 }d>0,\ \text{则必有 }\alpha\le \frac{c}{d}.
\]
因此，在强混合极限下不可能维持与临界线保持固定正距离的离线零点事件。
\end{theorem}

\begin{proof}
带外零点排除正是命题 \ref{prop:op-algebra-logsobolev-zero-band} 的内容。
另一方面，定理 \ref{thm:conclusion-critical-mellin-radial-double-rigidity} 已把 extension-preserving 且 unweighted 的可见 Mellin 切片唯一钉死在
\[
\Re(s)=\frac12.
\]
故一旦 \(\alpha\) 不退化，稳定零点事件只能停留在以临界线为中心的有界邻域中。
更精确地，若某稳定零点事件与临界线保持距离 \(d>0\)，则必有
\[
\alpha\le \frac{c}{d}.
\]
因此，在强混合极限下不可能维持与临界线保持固定正距离的离线零点事件。
\end{proof}

\begin{theorem}[连续横截唯一性与见证核的乘积闭包]\label{thm:conclusion-unique-continuous-transverse-and-witness-product-core}
在纯相位制度下，若允许外置连续模式轴以实现扩展保存，则任意使
\[
u\longmapsto \bigl(\operatorname{proj}_{\mathsf{Phase}}(u),r(u)\bigr)
\]
成为单射的连续记录轴 \(r\) 都必经由径向寄存器
\[
\varrho(u):=\log|u|
\]
因子化；
更强地，由定理 \ref{thm:conclusion-continuous-transverse-category-collapse}，任意有限个连续横截寄存器的联合仍塌缩为这一条连续轴。

另一方面，扩展保存记录轴的规范初对象
\[
r_0=r_{\mathrm{fin}}\times r_{\widehat{\ZZ}}\times r_{\mathrm{wind}}\times r_{\mathrm{end}}
\]
已由定理 \ref{thm:xi-extension-preserving-initial-object} 与定理 \ref{thm:cdim-minimal-record-axis} 给出。
因此，除去实现中立的有限部封口 \(r_{\mathrm{fin}}\) 之外，其见证相关的不可约核心至少包含三条互不可替代的离散因素：
\[
r_{\widehat{\ZZ}},
\qquad
r_{\mathrm{wind}},
\qquad
r_{\mathrm{end}}.
\]
换言之，纯相位制度下恰有且仅有一条连续横截方向，但完整可审计闭包仍至少需要三条正交的见证因素，而该唯一连续方向不能替代其中任何一条。
\end{theorem}

\begin{proof}
连续横截的唯一性与有限乘积塌缩由推论 \ref{cor:xi-unique-continuous-transverse-register} 及定理 \ref{thm:conclusion-continuous-transverse-category-collapse} 给出。
规范初对象的乘积闭包则由定理 \ref{thm:xi-extension-preserving-initial-object} 与定理 \ref{thm:cdim-minimal-record-axis} 给出。
最后，由定理 \ref{thm:xi-null-orthogonal-trichotomy-resource-nonsubstitutable}，碰撞方向、虚向折叠与端点连续累积所对应的见证预算在体系内彼此正交、不可替代，故 \(r_{\widehat{\ZZ}},r_{\mathrm{wind}},r_{\mathrm{end}}\) 的三条因素不能被压缩为更少的可审计核。
\leanverified{paper\_conclusion\_unique\_continuous\_transverse\_and\_witness\_product\_core}
\end{proof}

\begin{corollary}[不存在二轴审计的普适完备化]\label{cor:conclusion-no-universal-two-axis-audit-completion}
设一个纯相位协议同时允许以下三类现象非平凡出现：
\[
\text{同像碰撞},\qquad
\text{虚向折叠},\qquad
\text{端点连续累积}.
\]
则任何试图把全部不可逆性外置到至多两条记录轴上的普适审计方案都不可能既 extension-preserving 又保持完备。
更精确地：
\begin{enumerate}
  \item 若其中一条轴连续，则它在因子意义下只能等价于 \(\varrho\)，从而仍至多剩下一条离散轴，不足以同时承载 \(r_{\widehat{\ZZ}},r_{\mathrm{wind}},r_{\mathrm{end}}\)；
  \item 若两条轴都离散，则至多能覆盖三条见证因素中的两条，剩余一条必然失去可核验承载。
\end{enumerate}
\end{corollary}

\begin{proof}
由定理 \ref{thm:conclusion-unique-continuous-transverse-and-witness-product-core}，见证相关的最小闭包已经包含三条互不可替代的离散因素，而任意连续横截又只能落在唯一径向轴 \(\varrho\) 上。
因此，一切“至多两轴”的方案都必然遗漏至少一条正交见证因素，遂不可能成为普适完备化。
\end{proof}

\begin{conjecture}[自对偶多尺度审计的非周期临界刚性]\label{conj:conclusion-aperiodic-critical-rigidity}
设一完成化体系同时满足下列四项纪律：
\begin{enumerate}
  \item 自对偶函数方程与临界线实值性；
  \item 尺度混合落在非周期制度中，例如满足命题 \ref{prop:xi-scale-mellin-no-nonzero-imag-period} 的假设，或至少含有由定理 \ref{thm:xi-two-base-deperiodization-no-imag-period} 所支配的不可公度双底子混合；
  \item 不可见部分的对数 Sobolev 常数在所研究尺度上保持非退化下界；
  \item 扩展保存只允许经过定理 \ref{thm:cdim-minimal-record-axis} 所给的规范记录轴闭包，而不再引入新的连续横截方向。
\end{enumerate}
则其稳定零点几何应满足两条刚性：
\begin{enumerate}
  \item 全部稳定零点被压缩到以 \(\Re(s)=\frac12\) 为中心、由混合常数控制的临界邻域中；在强混合极限下，该邻域进一步塌缩到临界线本身；
  \item 零点集合不允许任何非零纯虚周期。
\end{enumerate}
换言之，可审计零点几何应是一种“临界线支持而竖向非晶格化”的对象：
既非单底 Dirichlet 梳齿，也非任意离线漂移的谱云，而是由自对偶性、混合强度、不可公度尺度与最小记录轴闭包共同锁定的非周期临界刚体。
\end{conjecture}

\begin{remark}[结构论据]\label{rem:conclusion-aperiodic-critical-rigidity-evidence}
定理 \ref{thm:conclusion-two-incommensurable-bases-destroy-vertical-lattice} 与定理 \ref{thm:conclusion-selfdual-scale-mellin-does-not-force-vertical-periodicity} 说明：函数方程与临界线实值性并不自带竖向晶格，而不可公度尺度已足以摧毁该晶格；
定理 \ref{thm:conclusion-logsobolev-mixing-forces-critical-localization} 说明：一旦混合强度保持非退化，稳定零点事件就被迫停留在临界线邻域；
定理 \ref{thm:conclusion-unique-continuous-transverse-and-witness-product-core} 与推论 \ref{cor:conclusion-no-universal-two-axis-audit-completion} 则说明：记录轴闭包本身也是刚性的，不能借由附加低维连续横截来偷偷恢复另一类谱几何。
四者合并，便得到“非周期临界刚性”这一统一图景。
\end{remark}

\endinput
